\section{Exactness of the unit-cost solutions}
\label{app:proofs}
The following arguments concern exact arithmetic, co-first and the assumptions in Section~\ref{sec:methods}: fixed rivals, nonnegative independently selectable target gains, equal task weights, and unit cost per selected coordinate. The implementation applies the documented tolerance and separately validates strict-first. No argument implies an ordering between different aggregation rules.

\paragraph{Arithmetic mean.}
Write $D=nH-\sum_jx_{tj}$ and $a_j=y_j-x_{tj}\geq0$. If $D\leq0$, the empty set is optimal. Otherwise order gains $a_{(1)}\geq\cdots\geq a_{(n)}$. Every size-$k$ set has gain sum at most $\sum_{r=1}^k a_{(r)}$, since replacing any selected smaller gain by an unselected larger gain cannot decrease that sum. Thus if this prefix falls below $D$, no size-$k$ set wins. The first prefix reaching $D$ is itself feasible and matches this lower bound. If the total falls below $D$, no subset can win. This proves optimality and the infeasibility criterion.

\paragraph{Geometric mean.}
If $H=0$, every target is already co-first relative to that threshold. If $H>0$, any unchanged zero forces the product to zero; hence every zero coordinate must be selected and must have $y_j>0$. If some such coordinate cannot become positive, the target is infeasible. Otherwise take those forced upgrades and define the positive starting vector $v$. Success is equivalent to $\sum_j\log x'_{tj}\geq n\log H$. All remaining gains $\log y_j-\log v_j$ are nonnegative, and selecting a coordinate adds its log gain. Applying the arithmetic exchange argument to the residual deficit gives the exact additional count. Forced selections are necessary, so the total is minimal.

\paragraph{Odd median.}
For $n=2h-1$, the median is at least $H$ if and only if at least $h$ entries are at least $H$. Let $q$ entries already satisfy this, and let $r$ currently failing coordinates be able to reach $H$ after improvement. Every selection increases the satisfying count by at most one. At least $\max(0,h-q)$ upgrades are therefore necessary. When $r$ is at least that deficit, selecting any such number of eligible coordinates suffices, attaining the lower bound. Otherwise the target is infeasible. At full unit ceilings all entries are eligible, so at most $h$ selections suffice. For even $n$ we use the actual middle-pair convention and exhaustive subsets; an odd-median count formula would not establish exactness there.

\paragraph{Minimum.}
The minimum is at least $H$ exactly when every coordinate is at least $H$. Hence every initially failing coordinate must be selected; if any cannot reach $H$, no feasible subset exists. Selecting precisely all failing coordinates otherwise suffices and is minimal.

\paragraph{Borda.}
Pairwise points decompose the total across tasks and opponents. On selected task $j$, the target's gain against opponent $i$ is $d_{ij}$. Pairwise points sum to one, including ties, so the opponent loses exactly $d_{ij}$. Target gain against all opponents is $g_j$. Thus after selection $z$, target-minus-opponent gap equals
\begin{equation}
 B_t-B_i+\sum_j(g_j+d_{ij})z_j.
\end{equation}
All rivals are beaten or tied exactly when each of these gaps is nonnegative, which is precisely the feasible set of Equation~\ref{eq:borda}. Minimizing $\sum_jz_j$ over binary $z$ therefore gives the desired minimum count. Half-integral original and modified Borda scores imply that strict positive gaps require at least half a point. A rival can lose at most one point per task and at most $n$ overall. If $M$ is the largest original rival total, one rival retains at least $M-n$; no rival originally below $M-n$ can become the strongest after nonnegative target improvements. Restricting constraints to rivals originally at least $M-n$ is therefore safe. The delivered solvers require optimal status with zero relative gap or a proven infeasible status; an unresolved optimization would not establish either result.

\paragraph{Monotonicity.}
Mean, median, geometric mean and minimum are coordinatewise nondecreasing in the target row, with fixed rival aggregates. For Borda, increasing a target coordinate weakly increases the target's pairwise points and weakly decreases every rival's points against it; other pairwise comparisons stay fixed. Relative standing therefore cannot worsen for any rule. If a menu permits larger improvements coordinatewise, every previously successful selected subset remains successful. The optimal count cannot increase, including the convention $\infty$ for infeasibility. This is why raising a cap cannot create a genuinely harder instance, even though conditioning descriptive summaries on its changing reachable cohort can produce apparent reversals.
